\begin{definition}[Open Mapping]
   Let \( X  \) and \( Y  \) be metric spaces. Then \( T: D(T) \to Y \) with domain \( D(T) \subseteq  X  \) is called an \textbf{open mapping} if for every open set in \( D(T) \) the image is an open set in \( Y  \). 
\end{definition}

\begin{lemma}[Open Unit Ball]
   A bounded linear operator \( T \) from a Banach space \( X  \) onto a Banach space \(  Y  \) has the property that the image \( T({B}_{0}) \) of the open ball \( {B}_{0} = B(0;1) \subseteq  X   \) contains an open ball about \( 0 \in Y \).
\end{lemma}

\begin{theorem}[Open Mapping Theorem/Bounded Inverse Theorem]
   A bounded linear operator \( T  \) from a Banach space \( X  \) onto a Banach space \( Y  \) is an open mapping. Hence, if \( T  \) is bijective, \( T^{-1} \) is continuous and thus bounded.
\end{theorem}
